% !TEX root = ../main.tex

\section{Preliminaries}
\label{sec:preliminaries}

\ecoparagraph{Notations.}
    We represent scalar discrete random variables over $N$ categories as one-hot column vectors, with ${\VC = \{x \in \{0,1\}^N\colon \sum_{i=1}^N x_i = 1\}}$ denoting the set of such vectors. We further define the probability simplex $\Delta \vcentcolon= \{ z \in \RR^N\colon z \ge 0, \langle z, \vec{1} \rangle = 1 \}$ as the convex hull of $\VC$, where $\vec{1} \in \RR^N$ denotes the all-ones vector. 

\subsection{Diffusion Language Models}
\label{sec:discrete diffusion models}
    Diffusion Language Models (DLMs; \citealp{austin2021structured, campbell2022continuous, lou2024discrete}) generate samples from a data distribution $p^*(x_0)$ by learning to reverse a prescribed forward noising process. In this work, we focus on two tractable process families that are widely used in language modeling: the \textit{absorbing} process, which underlies MDLM, and the \textit{uniform} process, which underlies UDLM and Duo. \reviewchange{We defer the general CTMC formulation and its SEDD score-entropy instantiation to Appendix~\ref{app:details of considered processes}. The process-specific details for masked diffusion and uniform diffusion are expanded in Appendices~\ref{app:absorbing process} and~\ref{app:uniform process}, respectively.}

\subsubsection{Absorbing Process}
\label{sec:absorbing process mdlm}
    \paragraph{Forward process.}
    The absorbing process corrupts tokens by replacing them with the mask token $m$. Its terminal distribution is concentrated on $m$, and under the usual schedule $\alpha_t \in [0,1]$, the forward transition from a clean token $x_0$ is
    \[
        q_t(x_t \mid x_0)
        = \mathrm{Cat}\bigl(x_t; \alpha_t x_0 + (1-\alpha_t)m\bigr).
    \]
    Thus, a token is either still equal to $x_0$ or has already become masked.

    \paragraph{Training loss.}
    MDLM~\citep{sahoo2024simple, shi2024simplified} parameterizes the denoiser as $f\colon \VC \times [0,1]\to\Delta$ and trains it with the weighted masked prediction objective
    \begin{equation}
    \label{eq:mdlm loss}
    \begin{aligned}
    \LC_{\text{MDLM}}(f, p^*)
    \vcentcolon= -
    \EE_{p^*(x_0), t, q_t}
    \bigl[
        \lambda_{t} \langle  x_0, \log f(x_t, t) \rangle
    \bigr].
    \end{aligned}
    \end{equation}
    where $\lambda_t$ is a positive weighting function and the logarithm is applied element-wise.

    \paragraph{Sampling procedure.}
    The \textsc{subs} parameterization of MDLM encodes two useful properties: unmasked tokens are copied forward, and the model assigns zero probability to predicting the mask token as clean data. For $0\leq s<t\leq1$, the learned reverse transition is the absorbing posterior with the clean token replaced by the model prediction, define $f_t\vcentcolon=f(x_t,t)$:
    \begin{equation}        
    \!\!p^{f}_{s\mid t}(x_s\!\mid\! x_t)
    \!\!=\!\!
    \begin{cases}
    \mathrm{Cat}(x_s; x_t),\!\! & \!\!x_t\neq m,\\[2pt]
    \mathrm{Cat}\!\left(
    \!x_s;
    \!\!\begin{gathered}
    \dfrac{(1\!-\!\alpha_s) m \! +\! (\alpha_s\!-\!\alpha_t)f_t}{1-\alpha_t}
    \end{gathered}\!\!
    \right), \!\! & \!\!x_t=m .
    \end{cases} \nonumber
    \end{equation}
    Sampling starts from a fully masked sequence $x_1\sim\delta_m(\cdot)$ and repeatedly draws $x_s\sim p^{f}_{s\mid t}(\cdot\mid x_t)$ along a decreasing time grid. Once a token is revealed, it is carried over in later denoising steps. The masked-diffusion sampling process is visualized in Figure~\ref{fig:masked-diffusion-sampling-intuition}. The full MDLM formulation is given in Appendix~\ref{app:mdlm-formulation-theory}.

\subsubsection{Uniform Process}
\label{sec:uniform process duo}
    \paragraph{Forward process.}
    The uniform process corrupts tokens toward the uniform distribution $\frac{1}{N}\vec{1}$. With the same notation $\alpha_t$, its forward transition can be written compactly as
    \[
        q_t(x_t \mid x_0)
        = \mathrm{Cat}\bigl(x_t; \alpha_t x_0 + (1-\alpha_t)\tfrac{1}{N}\vec{1}\bigr).
    \]
    Unlike absorbing diffusion, an intermediate token can be changed again at later reverse steps. This makes the process self-correcting: generation must decide both \textit{which} token to move to and \textit{when} to move.

    \paragraph{Training loss.}
    UDLM~\citep{schiff2024discreteguidance} uses the clean-token predictor as MDLM. Duo~\citep{sahoo2025the} keeps the same uniform discrete process but evaluates the model on a Gaussian-relaxed input, which gives smoother gradients during training while recovering the hard uniform process in the low-temperature limit, which makes the following parameterization of the model $f\colon \Delta \times [0,1]\to\Delta$. We write the objective as
    \begin{equation}
    \label{eq:duo loss}
    \LC_{\text{Duo}}(f,p^*)
    \vcentcolon=
      \textstyle
    \EE_{p^*(x_0), t, q_t}
      \bigl[
        g_{\mathrm{Duo}}\bigl(x_t ,x_0,f(x_t,t)\bigr)
      \bigr].
    \end{equation}
    All details of $g_{\mathrm{Duo}}$ are deferred to Appendix~\ref{app:uniform process}.

    \paragraph{Sampling procedure.}
    Sampling starts from a sequence of tokens sampled independently from the uniform distribution, $x_1\sim\mathrm{Cat}(\vec{1}/N)$. Then the sequence is iteratively refined by the learned uniform reverse posterior. For $0\leq s<t\leq1$, define $\alpha_{t\mid s}\vcentcolon=\alpha_t/\alpha_s$ and $f_t\vcentcolon=f(x_t,t)$. The ancestral sampler uses
    % {\small
    \begin{gather}
    \label{eq:uniform reversed}
    p^f_{s\mid t}(x_s\mid x_t) =
        \mathrm{Cat}\!\Bigl(x_s;
        \frac{N\alpha_t\,x_t\odot f_t
        +(\alpha_{t\mid s}-\alpha_t)x_t}{N\alpha_t\langle x_t,f_t\rangle+1-\alpha_t} \nonumber \\[-1pt]
        +\frac{(\alpha_s-\alpha_t)f_t
        +(1-\alpha_{t\mid s})(1-\alpha_s)\vec{1}/N}{{N\alpha_t\langle x_t,f_t\rangle+1-\alpha_t}} \Bigr). \nonumber
    \end{gather}
    % }
    \noindent In the sampling procedure each reverse step can revise the whole sequence.
    The uniform-diffusion sampling process is visualized in Figure~\ref{fig:appendix process intuition}. The full UDLM/Duo formulation is given in Appendix~\ref{app:uniform process}.
    
    The Duo paper also introduces an additional greedy sampling procedure. In the experiments, we report both variants and denote ancestral sampling by superscript ($^{\mathrm{a}}$) and greedy sampling by ($^{\mathrm{g}}$).


\subsubsection{Unified Notations \& Sequence Level Loss}
\label{sec:unified notations sequence loss}
    \paragraph{Unified token-level objective.}
    SEDD~\eqref{eq:sedd loss}, MDLM~\eqref{eq:mdlm loss}, and Duo~\eqref{eq:duo loss} use different losses, but they share the same structure. We write $\LC \in \{\LC_{\text{SEDD}}, \LC_{\text{MDLM}}, \LC_{\text{Duo}}\}$ in order to generalize our distillation scheme. Each loss can be expressed by the integrand $g$ of respective loss, so for any distribution $p$ we define:
    \[
    \begin{aligned}
    \LC(f,p)
    &=
    \EE_{p(x_0), t, q_t}\bigl[
        g\bigl(x_t,x_0,f(x_t,t)\bigr)
    \bigr].
    \end{aligned}
    \]

    \paragraph{Sequence-level objective.}
    In practical applications, we are interested in generating sequences $x_0^{1:L}$ of length $L$. To avoid an exponentially large state space, DLMs usually apply the forward corruption independently across sequence positions~\citep{lou2024discrete}. Under this assumption, the forward distribution factorizes as
    \[
        q_t(x_t^{1:L} \mid x_0^{1:L})
        = \prod_{l=1}^L q_t(x_t^l \mid x_0^l).
    \]
    The corresponding sequence-level diffusion loss becomes
    \begin{equation}
    \label{eq:sequence level loss}
    \begin{aligned}
    \LC^{1:L}(f,p^*)
    &=
    \EE_{p^*(x_0), t, q_t}\biggl[
        \sum_{l=1}^L
        g\bigl(
        x_t^l,x_0^l,f^l(x_t^{1:L},t)
        \bigr) \nonumber
    \biggr].
    \end{aligned}
    \end{equation}


\subsection{Inverse Distillation}
\label{sec:inverse distillation}

\ecoparagraph{Main goal.}
    The goal of inverse distillation is to directly train a student distribution $p_\theta(x_0)$, represented by a fast generator, so that it approximates the true data distribution $p^*(x_0)$. We use the pretrained diffusion model $f^*$, which was trained on samples from $p^*$. The good $p_\theta$ should be one for which the same diffusion training procedure would recover $f^*$.

\ecoparagraph{Diffusion training.}
    In usual diffusion training problem: given samples from $p^*$, we train the diffusion model by solving
    \begin{equation}
        f^* = \arg\min\nolimits_{f} \LC(f, p^*).
        \label{eq:teacher_objective_discr}
    \end{equation}
    Thus, the pretrained model $f^*$ is the diffusion model whose predictions are optimal for the real data distribution, we call this model the \textit{teacher}.

\ecoparagraph{Inverse viewpoint.}
    Inverse distillation reverses the roles of the known and unknown objects. We are given the pretrained diffusion model $f^*$ and seek a student distribution $p_\theta(x_0)$, such that training the same diffusion objective on samples from $p_\theta$ would recover the teacher:
    \begin{equation}
        f^* = \arg\min\nolimits_{f} \LC(f, p_{\theta}).
        \label{eq:inverse_optimality_condition}
    \end{equation}
    In words, diffusion training asks which diffusion model fits a distribution, while inverse distillation asks which distribution would make the fixed pretrained diffusion model optimal.

\ecoparagraph{Inverse distillation loss.}
    The optimality condition in~\eqref{eq:inverse_optimality_condition} is not directly feasible as a training objective, since it requires optimizing parameters $\theta$ through an $\argmin$. Following prior work on inverse distillation~\citep{gushchin2025inverse,kornilov2025universal, selikhanovych2025one}, we use tractable inverse loss
    \begin{gather}
    \begin{aligned}
    \LC_{\text{inv}}(\theta)
    \vcentcolon=
        \LC(f^*, p_{\theta})-
        \min\nolimits_{f} \LC(f, p_{\theta}).
    \end{aligned}
    \label{eq:inverse optimization problem}
    \end{gather}
    The inner minimization defines a \emph{fake} diffusion model trained on samples from $p_\theta$. The term $\LC(f^*,p_\theta)$ measures the teacher denoiser's performance on samples from $p_\theta$, while $\min_f \LC(f,p_\theta)$ is the best achievable denoising loss for a diffusion model trained on $p_\theta$. The generator is therefore optimized by minimizing this gap. By construction, $\LC_{\text{inv}}(\theta)\geq0$, and if $p_\theta=p^*$ then $\LC_{\text{inv}}(\theta)=0$. The converse is not guaranteed for this general objective.

\ecoparagraph{Connection to this work.}
    In \emph{continuous} space, inverse distillation was extensively explored by prior works~\citep{gushchin2025inverse,kornilov2025universal, selikhanovych2025one}. Our main goal is to extend this idea to the \emph{discrete} domain. However, it presents several challenges, stemming from the inherent characteristics of discrete spaces. These include \underline{theoretical concerns} regarding the validity of the training objective, due to the non-uniqueness of its minimizers, as well as \underline{practical difficulties} related to the non-differentiability of discrete sampling.
